\section{Evaluator-gap-driven preference mining}
\label{sec:mining}

\subsection{An evaluation semantic space}

Feedback acquisition begins with a versioned inventory of evaluation knowledge,
not with a convenient labeled dataset. For each domain, we record dimension
definitions, human protocols, rubrics, published metrics, learned scorers,
LLM/VLM judges, and external analysis tools. Each source has an input contract,
provenance, known limitations, and an execution status. A coverage audit checks
the resulting dimension--source matrix. Every core dimension needs executable
evidence; conflict mining ordinarily needs two sources with different mechanisms.
Several scorers sharing a model or training labels do not establish independent
coverage.

The unit of comparison is an \emph{output signal}, rather than an evaluator name.
We freeze a many-to-many mapping $M_{jq}$ from signal $j$ to dimension $q$
using its definition and measurement procedure before reading preference labels.
A mixed score may connect to several dimensions, but remains one shared signal.
Uninterpretable outputs are not assigned dimensions from post hoc correlations.
Existing executable metrics can enter the candidate tool library. Metrics
reserved for final independent assessment are frozen separately and never enter
tools, routing, critique, or reward.

\subsection{Semantic disagreement and structured annotation}

Raw scores on incompatible scales are not directly subtracted. For a candidate
pair $r=(x,y_a,y_b)$, each signal supplies a direction
$u_j(r)\in\{-1,0,+1,\bot\}$, with a normalized margin, repeated-run variance,
confidence, and availability. Zero denotes a tie and $\bot$ unavailable or
abstaining evidence. A small random anchor estimates reliability and human
disagreement before targeted sampling. For the reliable, nonabstaining signal
pairs $\mathcal{P}_q(r)$ mapped to the same dimension, a conflict score is
\begin{equation}
 c_q(r)=\frac{\sum_{(j,k)\in\mathcal{P}_q(r)}
             w_{jkq}\,\ind[u_j(r)\ne u_k(r)]}
            {\sum_{(j,k)\in\mathcal{P}_q(r)}w_{jkq}},
 \label{eq:conflict}
\end{equation}
where nonnegative weights reflect anchor reliability and cap the influence of
dependent sources. An empty denominator gives no conflict estimate.
This formulation leaves the calibration and dependence rule to be fixed before
the sampling comparison; disagreement is not itself a correctness label.

\begin{figure}[t]
\centering
\begin{tikzpicture}[line cap=round,line join=round]
  \path[use as bounding box] (0,0) rectangle (13.7,5.65);
  \pPanel{a}{Signal--dimension map}{0}{5.43}
  \pPanel{b}{Interpret disagreement}{6.03}{5.43}
  \pPanel{c}{Ask for evidence}{10.10}{5.43}

  \fill[pInk] (0,4.67) rectangle (5.75,5.10);
  \node[pSmall,text=white,anchor=west] at (.07,4.90) {Signal};
  \foreach \x/\t in {2.35/Content,3.32/Layout,4.29/Style,5.26/Overall}{
    \node[pSmall,text=white,align=center] at (\x,4.90) {\t};
  }
  \foreach \y/\lab in {
    4.38/OCR or text,
    3.95/CLIP (mixed),
    3.52/Position,
    3.09/SSIM (mixed),
    2.66/VLM content,
    2.23/VLM layout,
    1.80/VLM style,
    1.37/VLM overall}{
    \node[pSmall,anchor=west,text=pInk] at (.02,\y) {\lab};
    \foreach \x in {2.35,3.32,4.29,5.26}{
      \draw[pLine,line width=.35pt,fill=pPaper]
        (\x-.40,\y-.195) rectangle (\x+.40,\y+.195);
      \node[pSmall,text=pLine!90!black] at (\x,\y) {--};
    }
  }
  \foreach \x/\y in {2.35/4.38,2.35/3.95,4.29/3.95}{
    \fill[pBlue] (\x-.40,\y-.195) rectangle (\x+.40,\y+.195);
  }
  \foreach \x/\y in {3.32/3.52,3.32/3.09,4.29/3.09,2.35/2.66,3.32/2.23,4.29/1.80,5.26/1.37}{
    \fill[pRed] (\x-.40,\y-.195) rectangle (\x+.40,\y+.195);
  }
  \foreach \x/\y in {2.35/3.95,4.29/3.95,3.32/3.09,4.29/3.09}{
    \draw[pattern=north east lines,pattern color=white!36,draw=pLine]
      (\x-.40,\y-.195) rectangle (\x+.40,\y+.195);
  }
  \foreach \x/\y in {2.35/4.38,2.35/3.95,4.29/3.95}{
    \node[pTitle,text=white,fill=pBlue,inner sep=1pt] at (\x,\y) {A};
  }
  \foreach \x/\y in {3.32/3.52,3.32/3.09,4.29/3.09,2.35/2.66,3.32/2.23,4.29/1.80,5.26/1.37}{
    \node[pTitle,text=white,fill=pRed,inner sep=1pt] at (\x,\y) {B};
  }
  \node[pSmall,anchor=north west] at (.02,.94)
    {A/B: preference; --: not mapped.\\Hatching: one shared signal, not two votes.};

  \draw[pRule] (5.93,.58) -- (5.93,5.11);
  \foreach \y/\n in {4.67/1,3.18/2,1.69/3}{
    \node[pDot,draw=pLine,text=pMuted,minimum size=.35cm,
      font=\sffamily\fontsize{7}{8}\selectfont] at (6.30,\y) {\n};
  }
  \node[pLabel,anchor=west] at (6.57,4.67) {Within one dimension};
  \pCandidate{A}{pBlue}{6.52}{4.13}
  \node[pSmall,anchor=west] at (6.79,4.13) {OCR};
  \draw[pArrow] (7.47,4.13) -- (8.03,4.13);
  \draw[pArrow] (8.03,4.25) -- (7.47,4.25);
  \pCandidate{B}{pRed}{8.34}{4.13}
  \node[pSmall,anchor=west] at (8.61,4.13) {VLM};
  \node[pSmall,anchor=west] at (6.23,3.66) {Conflicting content signals};
  \draw[pRule] (6.23,3.43) -- (9.76,3.43);

  \node[pLabel,anchor=west] at (6.57,3.18) {Overall--dimension};
  \node[pText,anchor=west] at (6.23,2.64)
    {Overall \textcolor{pRed}{\textbf{B}} / content \textcolor{pBlue}{\textbf{A}}};
  \node[pSmall,anchor=west] at (6.23,2.18) {Check priorities and fusion};
  \draw[pRule] (6.23,1.96) -- (9.76,1.96);

  \node[pLabel,anchor=west] at (6.57,1.69) {Across dimensions};
  \node[pText,anchor=west] at (6.23,1.18)
    {Content \textcolor{pBlue}{\textbf{A}} / layout \textcolor{pRed}{\textbf{B}}};
  \node[pSmall,anchor=west] at (6.23,.72) {Trade-off, not contradiction};

  \node[pText,anchor=west] at (10.27,4.67) {$H_0$\quad Overall preference};
  \fill[pGold!9] (10.21,3.42) rectangle (13.7,4.27);
  \fill[pGold] (10.21,3.42) rectangle (10.28,4.27);
  \node[pText,anchor=west] at (10.40,4.04) {$\boldsymbol{H_1}$\quad Dimensional};
  \node[pSmall,text=pGold,anchor=west] at (10.40,3.67) {Core supervision};
  \node[pText,anchor=west] at (10.27,2.99) {$H_2$\quad Error + location};
  \node[pText,anchor=west] at (10.27,2.33) {$H_3$\quad Optional rationale};
  \draw[pEvidence] (11.96,2.00) -- (11.96,1.52);
  \node[pCard,draw=pTeal,fill=pTeal!5,align=center,text width=2.96cm,minimum height=.83cm]
    at (11.96,1.08) {\textbf{Evidence table $Z$}\\
    {\fontsize{7}{8.4}\selectfont\textcolor{pMuted}{H + signals + traces}}};
  \draw[pDash] (9.82,3.82) -- (10.16,3.82);
  \node[pSmall,align=center] at (6.85,.13)
    {Random reliability anchor $\rightarrow$ frozen acquisition policy $\rightarrow$ targeted conflicts; all annotation costs count};
\end{tikzpicture}

\caption{From evaluator signals to informative feedback. Panel (a) is a
synthetic webpage example of the many-to-many signal--dimension mapping; letters
indicate candidate preference, not scores. Hatched cells share one mixed signal.
Panel (b) distinguishes within-dimension contradiction, overall--dimension
disagreement, and a trade-off between different dimensions. Panel (c) connects
these strata to $H_0$--$H_3$ and the adaptation evidence table $Z$.
The matrix is illustrative and does not report measured scorer behavior.}
\label{fig:mining}
\end{figure}

We sample three separately recorded semantic strata
(Figure~\ref{fig:mining}b). \emph{Within-dimension disagreement} occurs when
reliable signals disagree about the same quality criterion.
\emph{Overall--dimension disagreement} indicates a possible mismatch in
priorities or aggregation. \emph{Cross-dimensional trade-offs} identify cases
where different criteria favor different candidates; they are not contradictions
between measurements of the same construct. Model differences, repeated-run
instability, and small margins further describe each stratum without replacing
its semantic definition. Candidate groups cover generators, decoding choices,
quality ranges, and failure types. After anchor calibration, the remaining pool's
labels stay hidden until the sampling policy is fixed. Only evolution data are conflict-sampled; final
evaluation uses prespecified random or stratified sampling.

\subsection{The adaptation evidence table}

Annotation targets overall and dimensional preferences, with localized
diagnostics on a conflict or failure subset. All candidates derived from one
input remain grouped. The adaptation evidence table $Z(r,q)$ joins artifact
provenance, human judgments and agreement, every source's direction and
uncertainty, evaluator decisions, selected routes, tool evidence, and execution
cost. This is the interface from preference mining to ERA.
Existing rankings, pointwise ratings, and error annotations require validated
adapters and explicit provenance; input/reference pairs alone provide no human
supervision. Appendix~\ref{app:annotation} specifies the feedback and readiness
audits. Neither model-generated explanations nor labels derived from a different
quality construct are promoted to human $H_1$ or $H_2$.
